\begin{abstract}
The growing photorealism of 3D generative models has made localized manipulation of CT volumes
(injecting synthetic lesions or erasing genuine findings) a credible clinical threat that
existing detectors fail to contain: supervised classifiers overfit to the generative architecture
seen at training, unsupervised approaches are structurally blind to removal attacks, and all
current methods produce opaque binary decisions incompatible with clinical transparency requirements.

We propose \textit{HEXMIL} (Hierarchical EXplainable Multiple Instance Learning), a weakly-supervised
framework that decomposes each CT volume into a three-level hierarchy of patches, axial slices,
and the full volume, applying independent gated attention at each level and training exclusively
on binary volume-level labels, without any pixel-level spatial supervision.

The product of the two learned attention maps assembles, as a direct by-product of classification,
a full-resolution 3D attention volume that spatially delineates the manipulated sub-volume: this
map is \emph{ante-hoc}, meaning it is the exact aggregation mechanism driving the binary prediction
and not a post-hoc saliency approximation, providing a faithfulness guarantee by construction.

On M3DSynth, spanning pix2pix, CycleGAN, and Latent Diffusion manipulations, HEXMIL achieves
superior in-domain classification and stronger out-of-distribution generalisation than existing
supervised and unsupervised detectors, while its emergent 3D attention maps outperform Grad-CAM
and GradCAM++ in pointing-game score under identical, zero-localization supervision.
\end{abstract}
